\section{Preliminaries: integration systems and the Daniell extension}\label{chap:preliminaries}

Section~\ref{chap:preliminaries} is a preliminary section whose purpose is to fix
the definitions, notation and statements used below.
Its content follows \cite{Yamagami}, and within it the construction of the
Daniell extension follows the procedure of \cite{BM}, as that book itself states.
\textbf{Nothing in Section~\ref{chap:preliminaries} is new.}

To keep the length down, Section~\ref{chap:preliminaries} states as a rule only
definitions and assertions, and we refer to the literature above for proofs.
There are two exceptions.

\begin{itemize}
  \item Lemma~\ref{lem:series-rep-equiv} (equivalence of the definition with
    dominating functions and the one without).
    The form of series representation that this paper adopts, following
    \cite{Yamagami}, differs on its face from the form adopted in \cite{BM}, so
    we give a proof as a bridge between the two.
  \item Lemma~\ref{lem:daniell-ext-is-complete} (Daniell closedness of the
    Daniell extension).
    It is used essentially in the main theorem of
    Section~\ref{chap:functoriality}, so we give a proof.
\end{itemize}


\subsection{Notation}

\begin{definition}
  When a sequence $\{a_n\}\subset\mathbb{R}$ is monotonically nonincreasing and
  $\lim_{n\to\infty}a_n$ converges to a real number $a$, we write
  \[
    a_n\downarrow a .
  \]
  Likewise, when a sequence $\{a_n\}\subset\mathbb{R}$ is monotonically
  nondecreasing and $\lim_{n\to\infty}a_n$ converges to a real number $a$, we write
  \[
    a_n\uparrow a .
  \]
\end{definition}

\begin{definition}
  Let $X$ be a set and $A$ a subset of $X$.
  The function
  \[
    1_A:X \to\mathbb{R};x\mapsto\begin{cases} 1 & x\in A \\ 0 & x\notin A \end{cases}
  \]
  is called the \textbf{indicator function} of the subset $A$\footnote{
    If $A=\emptyset$ then $x\notin\emptyset$ for every $x\in X$, so
    $1_\emptyset$ is the function constantly equal to $0$.
  }.
\end{definition}

\begin{definition}
  Let $X$ be a set.
  We define an order on the family $\mathbb{R}^X$ of functions on $X$ by
  \[
    f\le g \iff \forall x\in X  \; f(x)\le g(x)
  \]
  for all $f,g\in\mathbb{R}^X$.
\end{definition}

\begin{definition}
  We define operations $\land, \lor$ on the reals by
  \begin{align*}
    a\land b&:=\min\{a,b\}\\
    a\lor b&:=\max\{a,b\} ,
  \end{align*}
  and, for a set $X$ and $f,g\in\mathbb{R}^X$, we define the functions
  $f\land g$ and $f\lor g$ by
  \begin{align*}
    (f\land g)(x)&:=f(x)\land g(x) \\
    (f\lor g)(x)&:=f(x)\lor g(x) .
  \end{align*}
\end{definition}

\subsection{Riesz spaces and Riesz systems}

\begin{definition}
  Let $X$ be a set and let $L\subset\mathbb{R}^X$ be an
  $\mathbb{R}$-vector space\footnote{
    We allow $X=\emptyset$. In that case $\mathbb{R}^X=\{\emptyset\}$ is the
    zero-dimensional vector space, and $(\emptyset,\{0\},0)$ is the only
    integration system on the empty set.
    The definitions and assertions of this paper remain meaningful in this case,
    and they can be checked directly.
  }.
  We say that $L$ is a \textbf{Riesz space} on $X$\footnote{
    What this paper calls a Riesz space is called a vector lattice in some books.
    We use ``Riesz space'' throughout.
  } if
  \begin{align*}
    f\land g &\in L \\
    f\lor g &\in L
  \end{align*}
  holds for all $f,g\in L$.
\end{definition}

\begin{remark}[on terminology]\label{rem:riesz-terminology}
  The definition above is for subspaces of $\mathbb{R}^X$, and it requires
  closure under the lattice operations with respect to the pointwise order.
  Such an $L$ is also a Riesz space in the abstract sense (an ordered vector
  space that is a lattice), with the pointwise order as its order.
  That is, a ``Riesz space on $X$'' in this paper is a special case of an
  abstract Riesz space.

  In this paper, \textbf{whenever we attach ``on $X$'' we always mean a
  sublattice of $\mathbb{R}^X$}.
  Riesz spaces in the abstract sense appear only in
  Appendix~\ref{sec:forget}, where we call them \textbf{abstract Riesz spaces}
  to keep them apart.

  This distinction is not a formality, because not every abstract Riesz space can
  be realized as a Riesz space on some $X$.
  Indeed, if $L$ is a sublattice of $\mathbb{R}^X$ then evaluation
  $\mathrm{ev}_x$ at a point gives a Riesz homomorphism from $L$ to $\mathbb{R}$,
  and these separate the elements of $L$.
  But by Proposition~\ref{prop:spec-collapses} there is no nonzero Riesz
  homomorphism $L\to\mathbb{R}$ on $L^1[0,1]$.
  Hence $L^1[0,1]$ cannot be realized as a Riesz space on any set $X$.
\end{remark}

\begin{definition}
  Let $X$ be a set and let $L$ be a Riesz space on $X$.
  The pair
  \[
    (X , L)
  \]
  is called a \textbf{Riesz system} on the set $X$.
  The set $X$ in a Riesz system $(X,L)$ is called the \textbf{base space} of the
  Riesz system $(X,L)$.
\end{definition}

\begin{definition}
  Let $(X,L)$ and $(Y,M)$ be Riesz systems on sets $X$ and $Y$
  respectively.
  A \textbf{morphism of Riesz systems}
  \[
    F:(X,L)\to(Y,M)
  \]
  is a map $\underline{F}:Y\to X$ such that
  \[
    F^*f:=f\circ\underline{F}\in M
  \]
  holds for every $f\in L$.\footnote{Strictly, a morphism is the triple
  consisting of the domain, the codomain and the map of base spaces. Below we fix
  the domain and the codomain and present a morphism by the map of base spaces.}
\end{definition}

\begin{lemma}
  \begin{enumerate}
    \item For every Riesz system $(X,L)$ on a set $X$, the identity map
      $\mathrm{id}_X=\underline{\mathrm{id}_X}:X\to X$ induces a morphism of
      Riesz systems $\mathrm{id}_{(X,L)}:(X,L)\to (X,L)$.
    \item For any two morphisms of Riesz systems $F:(X,L)\to(Y,M)$ and
      $G:(Y,M)\to(Z,N)$, the map $\underline{F}\circ\underline{G}:Z\to X$ induces
      a morphism of Riesz systems $G\circ F:(X,L)\to(Z,N)$.
  \end{enumerate}
\end{lemma}
\begin{proof}
  \begin{enumerate}
    \item Immediate, since $f\circ \operatorname{id}_X=f\in L$ for every $f\in L$.
    \item For every $f\in L$ we have $F^*f\in M$ by hypothesis. Hence, again by
      hypothesis,
      $(G\circ F)^*(f)=f\circ(\underline{F}\circ\underline{G})=(f\circ \underline{F})\circ\underline{G}\in N$,
      which is the assertion.
  \end{enumerate}
\end{proof}

\begin{definition}
  The category whose objects are Riesz systems and whose morphisms are morphisms
  of Riesz systems is written
  \[
    \mathbf{RieszSys}
  \]
  and called the \textbf{category of Riesz systems}.
\end{definition}

\subsection{Integrals and integration systems}

\begin{definition}
  Let $X$ be a set and let $(X,L)$ be a Riesz system on $X$.
  We say that an $\mathbb{R}$-linear map
  \[
    I:L\to\mathbb{R}
  \]
  is an \textbf{integral} on the Riesz system $(X,L)$ if the following two
  conditions hold.
  \begin{itemize}
    \item For every $f\in L$, if $f\ge0$ then $I(f)\ge0$.
    \item For every sequence of functions $\{f_n\}\subset L$, if
      $f_n(x)\downarrow0$ for all $x\in X$ then $\lim_{n\to\infty}I(f_n)=0$.
  \end{itemize}
\end{definition}

\begin{definition}
  Let $(X,L)$ be a Riesz system on a set $X$ and let $I:L\to\mathbb{R}$
  be an integral on $(X,L)$.
  The triple
  \[
    (X,L,I)
  \]
  is called an \textbf{integration system} on the set $X$.
\end{definition}

\begin{example}[counting]\label{ex:counting}
  Let $X=\mathbb{N}$ and let $L$ be the set of all functions that vanish outside
  finitely many $n$.
  Clearly $L$ is a Riesz space on $X$.
  Putting
  \[
    I(f):=\sum_{n=1}^{\infty}f(n) ,
  \]
  the right-hand side is a finite sum and $I$ is an integral on $(X,L)$.
  Indeed, positivity is immediate, and if $f_n\downarrow0$ then the supports of
  the $f_n$ lie in the support of $f_1$, a finite set, so $I(f_n)\to0$.
\end{example}

\begin{example}[step functions and the Riemann integral]\label{ex:step-functions}
  Let $X=[0,1)$ and let $L$ be the set of all finite linear combinations of
  indicator functions of half-open intervals $[a,b)\subset[0,1)$, that is, the
  set of all step functions on $X$.
  Then $L$ is a Riesz space on $X$.
  Taking $I$ to be the Riemann integral
  \[
    I(f):=\int_0^1f(x)\,dx ,
  \]
  the triple $(X,L,I)$ is an integration system.

  Positivity is immediate. Continuity, that is, $I(f_n)\to0$ whenever
  $f_n(x)\downarrow0$ holds for every $x\in X$, follows from additivity of length
  over countable partitions of an interval.
  For a proof see \cite{BM}, \S2.
  The Daniell extension of this example is described in
  Example~\ref{ex:lebesgue}.
\end{example}

\begin{remark}\label{rem:why-half-open}
  There is a reason for taking the base space to be $[0,1)$ rather than $[0,1]$.
  If one takes $X=[0,1]$ and builds $L$ from indicator functions of half-open
  intervals, then every element of $L$ vanishes at $x=1$.
  Then $\sum_n\varphi_n(1)=0<+\infty$ for every series representation, so
  condition (3) of a series representation is imposed at $x=1$ always, and every
  element of $L^1$ satisfies $f(1)=0$.
  This differs from the set of all Lebesgue integrable functions on $[0,1]$.

  This phenomenon is a consequence of the fact that condition (3) of a series
  representation is shaped so as to constrain values only at the points where the
  sum of the dominating functions converges.
  If the base space has a point at which all elements of $L$ are forced to take
  the same value, then the Daniell extension becomes needlessly thin there.
\end{remark}

\begin{definition}
  Let $(X,L,I)$ and $(Y,M,J)$ be integration systems on sets $X$ and
  $Y$ respectively.
  A \textbf{morphism of integration systems}
  \[
    F:(X,L,I)\to(Y,M,J)
  \]
  is a morphism of Riesz systems $F:(X,L)\to(Y,M)$ such that
  \[
    J(F^*f)=I(f)
  \]
  holds for every $f\in L$.
\end{definition}

\begin{lemma}
  \begin{enumerate}
    \item For every integration system $(X,L,I)$ on a set $X$, the
      identity map $\operatorname{id}_X:X\to X$ induces a morphism of
      integration systems
      $\operatorname{id}_{(X,L,I)}:(X,L,I)\to (X,L,I)$.
    \item For any two morphisms of integration systems
      $F:(X,L,I)\to(Y,M,J)$ and $G:(Y,M,J)\to(Z,N,K)$, the map
      $\underline{F}\circ\underline{G}:Z\to X$ induces a morphism of integration
      systems $G\circ F:(X,L,I)\to(Z,N,K)$.
  \end{enumerate}
\end{lemma}
\begin{proof}
  We have already seen that the identity map induces a morphism of Riesz systems
  $\operatorname{id}_{(X,L)}$.
  Moreover $f\circ\underline{\operatorname{id}_X}=f\in L$ and
  $I(f\circ\underline{\operatorname{id}_X})=I(f)$ for every $f\in L$, so the
  identity map clearly induces a morphism of integration systems
  $\operatorname{id}_{(X,L,I)}:(X,L,I)\to (X,L,I)$.

  We have already seen that $\underline{F}\circ\underline{G}:Z\to X$ determines a
  morphism of Riesz systems from $(X,L)$ to $(Z,N)$.
  Moreover, by the hypothesis on morphisms of integration systems,
  \[
    K((G\circ F)^*(f))=K(G^*(F^*(f)))=J(F^*(f))=I(f) .
  \]
  Hence $\underline{F}\circ\underline{G}:Z\to X$ induces a morphism of
  integration systems $G\circ F:(X,L,I)\to(Z,N,K)$.
\end{proof}

\begin{definition}
  The category whose objects are integration systems and whose morphisms are
  morphisms of integration systems is written
  \[
    \mathbf{Int}
  \]
  and called the \textbf{category of integration systems}.
\end{definition}

\begin{remark}[on the direction of morphisms]\label{rem:opposite-convention}
  Note that the map of base spaces underlying a morphism
  $F:(X,L,I)\to(Y,M,J)$ is $\underline{F}:Y\to X$, so that the direction is
  reversed.
  In the language of measure theory, the condition
  $J(f\circ\underline{F})=I(f)$ says that $\underline{F}$ pushes $J$ forward to
  $I$, that is, that $\underline{F}$ is a measure-preserving map
  $(Y,J)\to(X,I)$.

  Hence \textbf{$\mathbf{Int}$ is the opposite of the usual category whose
  morphisms are measure-preserving maps}.
  This paper adopts this direction because pullback of functions is
  contravariant and we wish to treat the Daniell extension as a covariant
  functor.
  A reader who prefers to read everything in the usual convention should replace
  ``reflection'' below by ``coreflection'' and ``left adjoint'' by
  ``right adjoint''.
\end{remark}

\subsection{Integration systems satisfying the monotone convergence theorem}

\begin{definition}\label{def:satisfies-mct}
  We say that an integration system $(X,L,I)$ on a set $X$
  \textbf{satisfies the monotone convergence theorem} if, for every function
  $f:X\to\mathbb{R}$ and every monotonically nondecreasing sequence
  $\{f_n\}_{n=1}^{\infty}\subset L$,
  \[
    f_n(x)\uparrow f(x)\quad(\forall x\in X),\qquad\sup_n I(f_n)<+\infty
  \]
  implies
  \[
    f\in L,\qquad I(f)=\lim_{n\to\infty} I(f_n) .
  \]
\end{definition}

\noindent
This condition is a restatement in axiomatic form of the classical theorem of
Beppo Levi (1906), and \cite{CWS} calls a triple satisfying it \textbf{closed}
(its Definition 1.4.1).
The correspondence of terms with that book is collected in
Remark~\ref{rem:extended-real-collapse}.

\begin{definition}
  The full subcategory of the category $\mathbf{Int}$ of integration systems
  consisting of the integration systems that satisfy the monotone convergence
  theorem is written
  \[
    \mathbf{MInt}
  \]
  and called the \textbf{category of integration systems satisfying the monotone
  convergence theorem}.
\end{definition}

\subsection{The Daniell extension}

The content of this subsection follows \cite{Yamagami}. As that book states, the
facts that $L^1$ is a Riesz space and that $I$ extends to $L^1$ follow the
procedure of \cite{BM}.
\textbf{Nothing in this subsection is new; it is here to fix the definitions,
notation and assertions used in later subsections.}
Accordingly we state only definitions and assertions below, and refer to the
literature above for proofs.

\subsubsection{Definition of the Daniell extension}

\begin{definition}\label{def:series-representation}
  Let $(X,L,I)$ be an integration system on a set $X$.
  Given a function $f:X\to\mathbb{R}$, if there exist sequences of functions
  $\{f_n\}\subset L$ and $\{\varphi_n\}\subset L$ satisfying the following, then
  the pair of sequences $\{f_n\},\{\varphi_n\}$ is called a \textbf{series
  representation} of the function $f$, and a function admitting a series
  representation is said to be \textbf{integrable} with respect to the
  integration system $(X,L,I)$.
  \begin{enumerate}
    \item $|f_n|\le\varphi_n$ for every $n$;
    \item $\sum_{n=1}^\infty I(\varphi_n)<+\infty$;
    \item for every $x\in X$, if $\sum_{n=1}^\infty \varphi_n(x)<+\infty$ then
      $f(x)=\sum_{n=1}^\infty f_n(x)$.
  \end{enumerate}
  We write a series representation $\{f_n\},\{\varphi_n\}$ of a function $f$ as
  \[
    f\overset{\varphi_n}{\simeq}\sum f_n .
  \]
\end{definition}

\begin{definition}
  Let $(X,L,I)$ be an integration system on a set $X$.
  We denote the set of all functions integrable with respect to the integration
  system $(X,L,I)$ by
  \[
    L^1:=\{f\in\mathbb{R}^X \mid\text{$f$ is integrable with respect to the integration system $(X,L,I)$}\}
  \]
  \footnote{
    $L^1$ depends on $I$ as well, but in this paper we never put several
    integrals on one Riesz space $L$.
    We therefore write $L^1$, as something determined by $L$.
    Where this could be ambiguous we write it out, as in $(L^{\mathrm{M}})^1$.
    When referring to the object itself we use the functorial notation
    $\mathcal{D}(X,L,I)=(X,L^1,I^1)$.
  }.
\end{definition}

\medskip
Definition~\ref{def:series-representation} takes the form in which the
dominating functions $\varphi_n$ are made explicit, but this is equivalent to the
form without dominating functions adopted in \cite{BM}.
We will not use this below, but we state it in order to make the correspondence
with the literature clear.

\begin{lemma}[equivalence of the definitions with and without dominating functions]\label{lem:series-rep-equiv}
  Let $(X,L,I)$ be an integration system on a set $X$ and let
  $f:X\to\mathbb{R}$.
  Then the following are equivalent.
  \begin{enumerate}
    \item[(i)] $f$ is integrable in the sense of
      Definition~\ref{def:series-representation}.
      That is, there exist $\{f_n\},\{\varphi_n\}\subset L$ with
      $|f_n|\le\varphi_n$ and $\sum_n I(\varphi_n)<+\infty$ such that
      $f(x)=\sum_n f_n(x)$ holds at every $x$ with
      $\sum_n\varphi_n(x)<+\infty$.
    \item[(ii)] There exists $\{g_m\}\subset L$ with
      $\sum_m I(|g_m|)<+\infty$ such that $f(x)=\sum_m g_m(x)$ holds at every
      $x$ with $\sum_m|g_m(x)|<+\infty$.
  \end{enumerate}
\end{lemma}
\begin{proof}
  For $(ii)\Rightarrow(i)$ take $f_n:=g_n$ and $\varphi_n:=|g_n|$.
  Since $L$ is a Riesz space we have $|g_n|\in L$, and
  $\sum_n\varphi_n(x)<+\infty$ is literally $\sum_n|g_n(x)|<+\infty$.

  We show $(i)\Rightarrow(ii)$. Put $\psi_n:=\varphi_n-|f_n|\ge0$, so that
  $\psi_n\in L$.
  Define the sequence $\{g_m\}$ by
  \[
    f_1,\ \psi_1,\ -\psi_1,\ f_2,\ \psi_2,\ -\psi_2,\ f_3,\ \psi_3,\ -\psi_3,\ \dots .
  \]
  First, $I(|f_n|)\ge0$ because $I$ preserves order, and
  \[
    \sum_m I(|g_m|)
    =\sum_n\bigl(I(|f_n|)+2I(\psi_n)\bigr)
    =\sum_n\bigl(2I(\varphi_n)-I(|f_n|)\bigr)
    \le2\sum_n I(\varphi_n)<+\infty.
  \]
  Next, if $x$ satisfies $\sum_m|g_m(x)|<+\infty$, then
  \[
    \sum_m|g_m(x)|
    =\sum_n\bigl(|f_n(x)|+2\psi_n(x)\bigr)
    \ge\sum_n\bigl(|f_n(x)|+\psi_n(x)\bigr)
    =\sum_n\varphi_n(x)
  \]
  gives $\sum_n\varphi_n(x)<+\infty$. Hence $f(x)=\sum_n f_n(x)$ by condition
  (i).
  On the other hand $\sum_m g_m(x)$ converges absolutely, so rearrangement and
  regrouping are free, and $\psi_n$ cancels against $-\psi_n$, giving
  \[
    \sum_m g_m(x)=\sum_n f_n(x)=f(x) .
  \]
  Therefore (ii) holds.
\end{proof}

\subsubsection{Construction of the extended integral}

\begin{lemma}\label{lem:equiv-condition-for-riesz-space}
  Let $X$ be a set and let $L\subset\mathbb{R}^X$ be an
  $\mathbb{R}$-vector space.
  Then the condition that $|f|\in L$ for every $f\in L$ is equivalent to $L$
  being a Riesz space on $X$.
\end{lemma}

\begin{lemma}\label{lem:series-exp-is-abs-conv}
  Let $(X,L,I)$ be an integration system on a set $X$ and let
  $f\overset{\varphi_n}{\simeq}\sum f_n$.
  Then, for every $x\in X$, if $\sum_{n=1}^\infty \varphi_n(x)<+\infty$ then
  $f(x)=\sum_{n=1}^\infty f_n(x)$ converges absolutely.
\end{lemma}

\begin{theorem}\label{thm:l1-is-riesz-space}
  Let $(X,L,I)$ be an integration system on a set $X$.
  Then the family $L^1$ of functions integrable with respect to the integration
  system $(X,L,I)$ is a Riesz space on $X$.
\end{theorem}

\begin{lemma}\label{lem:integration-preserve-order}
  Let $(X,L,I)$ be an integration system on a set $X$.
  Then $I(f)\le I(g)$ for all $f,g\in L$ with $f\le g$.
\end{lemma}

\begin{lemma}\label{lem:triangle-ineq-for-integration}
  Let $(X,L,I)$ be an integration system on a set $X$.
  Then $|I(f)|\le I(|f|)$ for every $f\in L$.
\end{lemma}

\begin{lemma}\label{lem:abs-conv-for-int-of-series-rep}
  Let $(X,L,I)$ be an integration system on a set $X$.
  For every integrable function $f\in L^1$ with respect to the integration
  system $(X,L,I)$ and every series representation
  $f\overset{\varphi_n}{\simeq}\sum f_n$ of it, the sum
  $\sum_{n=1}^\infty I(f_n)$ converges absolutely.
\end{lemma}

\begin{lemma}[continuity of the integral]\label{lem:integration-is-continuous}
  Let $(X,L,I)$ be an integration system on a set $X$, let $\{f_n\}$ be
  a sequence of functions in $L$, and let $f\in L$.
  Then, if $f_n(x)\uparrow f(x)$ for every $x\in X$, we have
  $I(f_n)\uparrow I(f)$.
\end{lemma}

\begin{lemma}\label{lem:ext-int-of-zero-func}
  Let $(X,L,I)$ be an integration system on a set $X$.
  Then, given a series representation $0\overset{\eta_n}{\simeq}\sum h_n$ of the
  function constantly equal to $0$,
  \[
    \sum_{n=1}^\infty I(h_n)=0 .
  \]
\end{lemma}

\begin{lemma}\label{lem:ext-int-well-defined}
  Let $(X,L,I)$ be an integration system on a set $X$.
  Then, for any two series representations
  \begin{align*}
    f&\overset{\varphi_n}{\simeq}\sum f_n \\
    f&\overset{\psi_n}{\simeq}\sum g_n
  \end{align*}
  of an integrable function $f\in L^1$ with respect to the integration system
  $(X,L,I)$, we have $\sum_{n=1}^\infty I(f_n)=\sum_{n=1}^\infty I(g_n)$.
\end{lemma}

\begin{definition}
  Let $(X,L,I)$ be an integration system on a set $X$.
  For each integrable function $f\in L^1$ with respect to the integration system
  $(X,L,I)$, choose a series representation
  $f\overset{\varphi_n}{\simeq}\sum f_n$.
  Then the absolutely convergent sum
  (Lemma~\ref{lem:abs-conv-for-int-of-series-rep})
  \[
    \sum_{n=1}^\infty I(f_n)
  \]
  does not depend on the choice of series representation, by
  Lemma~\ref{lem:ext-int-well-defined}.
  We define this value to be $I^1(f)$.
\end{definition}

\begin{lemma}\label{lem:ext-int-is-linear}
  Let $(X,L,I)$ be an integration system on a set $X$.
  Then the map $I^1:L^1\to\mathbb{R}$ is an $\mathbb{R}$-linear map on the Riesz
  space $L^1$ on $X$.
\end{lemma}

\begin{lemma}[truncation estimate for a series representation of a positive function]\label{lem:positive-series-representation-truncation}
  Let $(X,L,I)$ be an integration system on a set $X$ and let
  $f\in L^1$ be an integrable function with $f\ge0$.
  Choose a series representation $f\overset{\varphi_n}{\simeq}\sum f_n$ of $f$
  and, for each $N\ge1$, put
  \[
    S_N:=\sum_{k=1}^Nf_k,\qquad h_N:=S_N\vee0,\qquad q_N:=S_N\wedge0 .
  \]
  Then
  \[
    I(|q_N|)\le\sum_{k=N+1}^\infty I(\varphi_k),
    \qquad
    |I^1(f)-I(h_N)|\le2\sum_{k=N+1}^\infty I(\varphi_k)
  \]
  hold. In particular $I(h_N)\longrightarrow I^1(f)$ as $N\to\infty$.
\end{lemma}

\begin{lemma}\label{lem:ext-int-is-positive}
  Let $(X,L,I)$ be an integration system on a set $X$ and let
  $f\in L^1$.
  Then $f\ge0$ implies $I^1(f)\ge0$.
\end{lemma}

\begin{lemma}\label{lem:positive-good-representation}
  Let $(X,L,I)$ be an integration system on a set $X$ and let
  $f\in L^1$.
  Then, if $f\ge0$, for every $\varepsilon>0$ there is a series representation
  $f\overset{\varphi_n}{\simeq}\sum f_n$ of $f$ satisfying
  \[
    \sum_{n=1}^\infty I(\varphi_n)\le I^1(f)+\varepsilon .
  \]
\end{lemma}

\begin{lemma}[flattening series representations into one sequence]\label{lem:flatten-series-representations}
  Let $(X,L,I)$ be an integration system on a set $X$.
  Suppose that for each $n\ge1$ a series representation
  $u_n\overset{\varphi_{n,k}}{\simeq}\sum_{k=1}^\infty g_{n,k}$ of $u_n\in L^1$
  is given, and that
  \[
    \sum_{n=1}^\infty\sum_{k=1}^\infty I(\varphi_{n,k})<+\infty .
  \]
  Put
  \[
    A:=\left\{x\in X\mathrel{}\middle|\mathrel{}\sum_{n=1}^\infty\sum_{k=1}^\infty\varphi_{n,k}(x)<+\infty\right\} .
  \]
  Then $\sum_{n=1}^\infty u_n(x)$ converges absolutely for every $x\in A$.

  Suppose further that a function $u:X\to\mathbb{R}$ satisfies
  $u(x)=\sum_{n=1}^\infty u_n(x)$ for $x\in A$.
  Choose a bijection $\sigma:\mathbb{N}\to\mathbb{N}\times\mathbb{N}$,
  $\sigma(m)=(n_m,k_m)$, and put $\widetilde{g}_m:=g_{n_m,k_m}$ and
  $\widetilde{\varphi}_m:=\varphi_{n_m,k_m}$.
  Then $u\overset{\widetilde{\varphi}_m}{\simeq}\sum_{m=1}^\infty\widetilde{g}_m$
  is a series representation of $u$.
  In particular $u\in L^1$ and
  \[
    I^1(u)=\sum_{n=1}^\infty I^1(u_n)
  \]
  holds, the right-hand side converging absolutely.
\end{lemma}

\begin{theorem}\label{thm:ext-int-is-integral}
  Let $(X,L,I)$ be an integration system on a set $X$.
  Then the map $I^1:L^1\to\mathbb{R}$ is an integral on the Riesz space $L^1$ on
  $X$.
\end{theorem}

\begin{lemma}\label{lem:daniell-ext-contains-source}
  Let $(X,L,I)$ be an integration system on a set $X$ and let $f\in L$.
  Then $f\in L^1$ and $I^1(f)=I(f)$.
\end{lemma}

\begin{definition}
  Let $(X,L,I)$ be an integration system on a set $X$.
  The integration system $(X,L^1,I^1)$ is called the \textbf{Daniell extension}
  of the integration system $(X,L,I)$.
\end{definition}

\begin{example}[the Lebesgue integral]\label{ex:lebesgue}
  Consider the integration system of Example~\ref{ex:step-functions}, that is,
  the integration system $(X,L,I)$ consisting of the step functions on $[0,1)$
  and the Riemann integral.
  Then its Daniell extension $(X,L^1,I^1)$ coincides with the set of all
  Lebesgue integrable functions on $[0,1)$ together with the Lebesgue integral.

  That is, the construction of Section~\ref{chap:preliminaries} recovers the
  Lebesgue integral.
  This is the origin of the definition of integrability by series
  representations; for details see \cite{Mikusinski} and \cite{BM}, \S2.
\end{example}

\begin{corollary}
  Let $(X,L,I)$ be an integration system on a set $X$ and let
  $(X,L^1,I^1)$ be its Daniell extension.
  Then the identity map $\operatorname{id}_X$ induces a morphism of integration
  systems $i_X:(X,L,I)\to(X,L^1,I^1)$.
\end{corollary}

\subsubsection{The Daniell extension satisfies the MCT}

\begin{lemma}[termwise integration of a nonnegative series]\label{lem:positive-series-flattening}
  Let $(X,L,I)$ be an integration system on a set $X$ and let
  $(X,L^1,I^1)$ be its Daniell extension.
  Suppose $\{h_k\}_{k=1}^{\infty}\subset L^1$ with $h_k\ge0$ satisfies
  $\sum_k I^1(h_k)<\infty$.
  Then, putting
  \[
    h(x):=\begin{cases}
      \sum_n h_n(x) & \text{if $\sum_n h_n(x)$ converges} \\
      0 & \text{otherwise,}
    \end{cases}
  \]
  we have $h\in L^1$ and $I^1(h)=\sum_k I^1(h_k)$.
\end{lemma}

\begin{lemma}\label{lem:daniell-ext-satisfies-mct}
  Let $(X,L,I)$ be an integration system on a set $X$.
  Then its Daniell extension $(X,L^1,I^1)$ satisfies the monotone convergence
  theorem (Definition~\ref{def:satisfies-mct}).
\end{lemma}

\subsection{The category of Daniell closed integration systems}

\begin{definition}\label{def:daniell-closed}
  We say that an integration system $(X,L,I)$ on a set $X$ is
  \textbf{Daniell closed} if
  \[
    L^1=L .
  \]
\end{definition}

\begin{remark}[on terminology]\label{rem:why-daniell-closed}
  This condition is often called ``completeness''.
  We assign a different word to it because
  \textbf{``complete'' occurs in this paper in three senses}:
  \begin{itemize}
    \item $L^1=L$ (the present definition);
    \item completeness with respect to the norm $\|\cdot\|_I$
      (Theorem~\ref{thm:l1-norm-complete});
    \item completion of a measure (Subsection~\ref{sec:why-complete},
      Remark~\ref{rem:separation-not-sufficient}).
  \end{itemize}
  These are three distinct conditions, and in this paper all three live in the
  same subsection.
  We therefore call the first one \textbf{Daniell closed}.
  The name follows \cite{CWS} and \cite{CFW}, which use the word ``closure'' of a
  Daniell space.

  \textbf{It also clashes with ``complete'' in other books.}
  \cite{CWS} uses ``complete'' for the sandwich condition ($f\le h\le g$
  together with $\ell(f)=\ell(g)$ implies that $h$ belongs as well), which is
  none of the three above but rather our
  Definition~\ref{def:mod-closed}.
  That book further uses ``complete'' for order-complete lattices and for the
  completion of a measure space.
  The correspondence of terms is collected in
  Remark~\ref{rem:extended-real-collapse}.

  In \textbf{the title}, however, \textbf{we follow custom and say
  ``completeness''}.
  Outside this paper that is the standard word for $L^1=L$, and the coexistence
  of the three senses is an internal matter of this paper.
  In the body and the abstract we use ``Daniell closed'' consistently.
\end{remark}

\begin{remark}\label{rem:complete-implies-mct}
  A Daniell closed integration system automatically satisfies the monotone
  convergence theorem.
  Indeed, the Daniell extension $(X,L^1,I^1)$ satisfies the monotone convergence
  theorem by Lemma~\ref{lem:daniell-ext-satisfies-mct}, and since $L^1=L$ and
  $I^1=I$ (Lemma~\ref{lem:daniell-ext-contains-source}), the system $(X,L,I)$
  itself satisfies the monotone convergence theorem.

  We may therefore define $\mathbf{DInt}$ as a full subcategory of
  $\mathbf{MInt}$.
  Where we list the monotone convergence theorem and Daniell closedness side by
  side below, it is in order to make the relation to $\mathbf{MInt}$ easy to see.
\end{remark}

\begin{definition}
  The full subcategory of the category $\mathbf{MInt}$ of integration systems
  satisfying the monotone convergence theorem consisting of the Daniell closed
  integration systems is written
  \[
    \mathbf{DInt}
  \]
  and called the \textbf{category of Daniell closed integration systems}.
\end{definition}

\begin{definition}\label{def:inclusion-functors}
  We write the inclusion functors among the three categories above as
  \begin{align*}
    \iota&:\mathbf{MInt}\to\mathbf{Int},\\
    \iota_{D}&:\mathbf{DInt}\to\mathbf{Int},\\
    \iota_M&:\mathbf{DInt}\to\mathbf{MInt}
  \end{align*}
  respectively. All of them are inclusions of full subcategories, and
  $\iota_{D}=\iota\circ\iota_M$.
\end{definition}

\begin{lemma}\label{lem:daniell-ext-is-complete}
  Let $(X,L,I)$ be an integration system on a set $X$ and let
  $(X,L^1,I^1)$ be its Daniell extension.
  Then the integration system $(X,L^1,I^1)$ is Daniell closed.
\end{lemma}

\noindent
This statement itself is stated as a corollary in \cite{Yamagami} and is not new.
However, since this paper uses it essentially in the main theorem of
Section~\ref{chap:functoriality}, in the form ``the codomain of the
factorization is Daniell closed'', we give a proof in this subsection only.
\begin{proof}
  We have already checked that the monotone convergence theorem holds, so it
  suffices to show $(L^1)^1=L^1$.
  Since $L^1\subset(L^1)^1$ has already been shown
  (Lemma~\ref{lem:daniell-ext-contains-source}), it suffices to show
  $(L^1)^1\subset L^1$.

  Let $f\in (L^1)^1$. That is, there are sequences of functions
  $\{\varphi_n\},\{f_n\}\subset L^1$ such that
  \begin{enumerate}
    \item $|f_n|\le\varphi_n$;
    \item $\sum_{n=1}^\infty I^1(\varphi_n)<+\infty$;
    \item if $\sum_{n=1}^\infty \varphi_n(x)<+\infty$ then
      $f(x)=\sum_{n=1}^\infty f_n(x)$.
  \end{enumerate}
  For each $n$ put
  \[
    f_n^+:=f_n\lor0,
    \qquad
    f_n^-:=(-f_n)\lor0 .
  \]
  Since $L^1$ is a Riesz space we have $f_n^+,f_n^-\in L^1$, and
  \[
    f_n=f_n^+-f_n^-,
    \qquad
    |f_n|=f_n^++f_n^- .
  \]
  Moreover, by positivity of $I^1$ and $|f_n|\le\varphi_n$,
  \[
    \begin{aligned}
      \sum_{n=1}^\infty I^1(f_n^+)
      +\sum_{n=1}^\infty I^1(f_n^-)
      &=\sum_{n=1}^\infty I^1(|f_n|)\\
      &\le\sum_{n=1}^\infty I^1(\varphi_n)
      <+\infty.
    \end{aligned}
  \]

  Now define functions $p,q:X\to\mathbb{R}$ by
  \[
    p(x):=
    \begin{cases}
      \sum_n f_n^+(x) & \text{if $\sum_n f_n^+(x)$ converges},\\
      0 & \text{otherwise}
    \end{cases}
  \]
  and
  \[
    q(x):=
    \begin{cases}
      \sum_n f_n^-(x) & \text{if $\sum_n f_n^-(x)$ converges},\\
      0 & \text{otherwise.}
    \end{cases}
  \]
  By Lemma~\ref{lem:positive-series-flattening} we have $p,q\in L^1$.

  Put
  \[
    A:=\left\{x\in X\mathrel{}\middle|\mathrel{}\sum_{n=1}^\infty\varphi_n(x)<+\infty\right\} .
  \]
  If $x\in A$ then
  \[
    \sum_{n=1}^\infty|f_n(x)|
    \le
    \sum_{n=1}^\infty\varphi_n(x)
    <+\infty ,
  \]
  so the two series defining $p(x)$ and $q(x)$ both converge. Hence, by the
  condition on the original series representation,
  \[
    f(x)=\sum_{n=1}^\infty f_n(x)=p(x)-q(x)
    \qquad(x\in A).
  \]
  Putting
  \[
    r:=f-(p-q) ,
  \]
  the function $r$ vanishes on $A$.

  Choose $\varepsilon_n>0$ for each $n$ so that
  \[
    \sum_{n=1}^\infty\varepsilon_n<+\infty .
  \]
  Since $\varphi_n\ge0$, by Lemma~\ref{lem:positive-good-representation} we may
  choose for each $n$ a series representation
  \[
    \varphi_n\overset{\psi_{n,k}}{\simeq}\sum_k g_{n,k}
  \]
  of $\varphi_n$ satisfying
  \[
    \sum_{k=1}^\infty I(\psi_{n,k})
    \le I^1(\varphi_n)+\varepsilon_n .
  \]
  Hence
  \[
    \sum_{n,k}I(\psi_{n,k})
    \le
    \sum_n I^1(\varphi_n)+\sum_n\varepsilon_n
    <+\infty.
  \]

  For each $n$,
  \[
    0\overset{\psi_{n,k}}{\simeq}\sum_k0
  \]
  is a series representation of the function constantly equal to $0$. Put
  \[
    B:=
    \left\{
      x\in X
      \mathrel{}\middle|\mathrel{}
      \sum_{n,k}\psi_{n,k}(x)<+\infty
    \right\} .
  \]
  If $x\in B$ then $\sum_k\psi_{n,k}(x)<+\infty$ for each $n$, so
  \[
    \varphi_n(x)=\sum_k g_{n,k}(x)
  \]
  and
  \[
    0\le\varphi_n(x)
    \le\sum_k|g_{n,k}(x)|
    \le\sum_k\psi_{n,k}(x).
  \]
  Hence
  \[
    \sum_n\varphi_n(x)
    \le\sum_{n,k}\psi_{n,k}(x)
    <+\infty.
  \]
  Therefore $x\in A$, and
  \[
    r(x)=0=\sum_n0
    \qquad(x\in B).
  \]
  Applying Lemma~\ref{lem:flatten-series-representations} with $u_n=0$ and
  $u=r$, we get $r\in L^1$.

  Since $p,q,r\in L^1$ and $L^1$ is a vector space,
  \[
    f=p-q+r\in L^1 .
  \]
  Therefore $(L^1)^1\subset L^1$, and $(L^1)^1=L^1$ is proved.
\end{proof}

\subsubsection{Splitting Daniell closedness into two closure conditions}

Daniell closedness $L^1=L$ is a single closure condition passing through the
extension by series representations.
In this subsubsection we show that this condition splits into two independent
conditions, \textbf{closure under the monotone convergence theorem} and
\textbf{closure under modification on null sets}
(Theorem~\ref{thm:complete-intrinsic}).
That neither condition can be dropped is seen in
Remark~\ref{rem:two-conditions-independent}.

The two conditions on the right are conditions that occur independently in
measure theory.
The monotone convergence theorem corresponds directly, and closure under
modification on null sets corresponds to completeness of the measure.
This decomposition is important for gauging the content of the main theorem of
Section~\ref{chap:functoriality}
(Remark~\ref{rem:main-theorem-not-tautology}).

\begin{definition}\label{def:null-set}
  Let $(X,L,I)$ be an integration system on a set $X$.
  We say that a subset $Z\subset X$ is an \textbf{$I$-null set} if there exists a
  sequence of functions $\{g_n\}$ in $L$ satisfying the following three
  conditions.
  \begin{enumerate}
    \item $g_n\ge0$ for every $n$;
    \item $\sum_{n=1}^{\infty}I(g_n)<+\infty$;
    \item $\sum_{n=1}^{\infty}g_n(x)=+\infty$ for every $x\in Z$.
  \end{enumerate}
\end{definition}

\noindent
This definition does not mention $L^1$.
Note that it is written using only the data given by $(X,L,I)$.
The set of all points at which condition (3) of a series representation
(Definition~\ref{def:series-representation}) does not constrain the value is
always an $I$-null set, as one sees by taking $\{g_n\}$ to be the dominating
functions $\{\varphi_n\}$.

\begin{lemma}\label{lem:null-sets-countable-union}
  A countable union of $I$-null sets is an $I$-null set.
\end{lemma}
\begin{proof}
  Let $Z_1,Z_2,\dots$ be $I$-null sets, choose for each $Z_k$ a sequence
  $\{g^{(k)}_n\}_n$ as in Definition~\ref{def:null-set}, and put
  $C_k:=\sum_nI(g^{(k)}_n)<+\infty$.
  Put $a_k:=2^{-k}/(1+C_k)>0$ and $\tilde{g}^{(k)}_n:=a_kg^{(k)}_n\in L$.
  Arranging the family $\{\tilde{g}^{(k)}_n\}_{k,n}$ into a single sequence,
  every term is nonnegative and
  \[
    \sum_{k,n}I(\tilde{g}^{(k)}_n)=\sum_ka_kC_k\le\sum_k2^{-k}=1<+\infty .
  \]
  Moreover, if $x\in\bigcup_kZ_k$ then $x\in Z_k$ for some $k$ and
  \[
    \sum_n\tilde{g}^{(k)}_n(x)=a_k\sum_ng^{(k)}_n(x)=+\infty .
  \]
  The remaining terms are nonnegative, so the sum of the rearranged sequence is
  $+\infty$ as well.
\end{proof}

\begin{lemma}\label{lem:vanish-off-null-set}
  Let $(X,L,I)$ be an integration system and let $Z\subset X$ be an $I$-null set.
  If $g\in L$ vanishes on $X\setminus Z$ then
  \[
    I(|g|)=0,
    \qquad\text{and in particular}\qquad
    I(g)=0 .
  \]
\end{lemma}
\begin{proof}
  Choose $\{g_n\}$ as in Definition~\ref{def:null-set} and put
  $C:=\sum_nI(g_n)<+\infty$.
  Take an arbitrary $\varepsilon>0$ and put
  \[
    u_N:=|g|\land\Bigl(\varepsilon\sum_{n\le N}g_n\Bigr) .
  \]
  Since $L$ is a Riesz space we have $u_N\in L$.

  We check that $u_N$ increases monotonically to $|g|$ pointwise.
  If $x\notin Z$ then $|g|(x)=0$, and the second term is nonnegative, so
  $u_N(x)=0=|g|(x)$.
  If $x\in Z$ then $\varepsilon\sum_{n\le N}g_n(x)\uparrow+\infty$, so
  $u_N(x)=|g|(x)$ for $N$ large enough, while before that
  $u_N(x)=\varepsilon\sum_{n\le N}g_n(x)$, which is nondecreasing.
  In either case $u_N(x)\uparrow|g|(x)$.

  Since $|g|\in L$, Lemma~\ref{lem:integration-is-continuous} gives
  $I(u_N)\uparrow I(|g|)$.
  On the other hand $u_N\le\varepsilon\sum_{n\le N}g_n$, so
  Lemma~\ref{lem:integration-preserve-order} gives
  \[
    I(u_N)\le\varepsilon\sum_{n\le N}I(g_n)\le\varepsilon C.
  \]
  Hence $I(|g|)\le\varepsilon C$. As $\varepsilon>0$ was arbitrary, $I(|g|)=0$,
  and $|I(g)|\le I(|g|)=0$ by
  Lemma~\ref{lem:triangle-ineq-for-integration}.
\end{proof}

\begin{definition}\label{def:mod-closed}
  We say that an integration system $(X,L,I)$ on a set $X$ is
  \textbf{closed under modification on null sets} if
  \[
    h|_{X\setminus Z}=f|_{X\setminus Z}
    \quad\Longrightarrow\quad
    h\in L
  \]
  holds for every $f\in L$, every $I$-null set $Z\subset X$ and every function
  $h:X\to\mathbb{R}$.
\end{definition}

\noindent
On $Z$ the function $h$ is subject to no constraint at all; only that it is
real-valued.

\begin{theorem}[decomposition of Daniell closedness]\label{thm:complete-intrinsic}
  Let $(X,L,I)$ be an integration system on a set $X$. The following are
  equivalent.
  \begin{enumerate}
    \item[(i)] $(X,L,I)$ is Daniell closed, that is, $L^1=L$.
    \item[(ii)] $(X,L,I)$ satisfies the monotone convergence theorem and is
      closed under modification on null sets.
  \end{enumerate}
\end{theorem}
\begin{proof}
  \textbf{(i)$\Rightarrow$(ii)}.
  The monotone convergence theorem was seen in
  Remark~\ref{rem:complete-implies-mct}.
  We show closure under modification on null sets.
  Let $f\in L$, let $Z$ be an $I$-null set, and let $h:X\to\mathbb{R}$ agree with
  $f$ on $X\setminus Z$.
  Choose $\{g_n\}$ as in Definition~\ref{def:null-set} and put
  \[
    f_1:=f,\quad f_n:=0\ (n\ge2),
    \qquad
    \varphi_1:=|f|,\quad \varphi_{n+1}:=g_n\ (n\ge1) .
  \]
  Then $|f_n|\le\varphi_n$ is clear, and
  \[
    \sum_{n}I(\varphi_n)=I(|f|)+\sum_nI(g_n)<+\infty .
  \]
  Moreover, if $\sum_n\varphi_n(x)<+\infty$ then $\sum_ng_n(x)<+\infty$, so
  $x\notin Z$ and therefore $h(x)=f(x)=\sum_nf_n(x)$.
  Hence $\{f_n\},\{\varphi_n\}$ is a series representation of $h$, and
  $h\in L^1=L$.

  \textbf{(ii)$\Rightarrow$(i)}.
  Since $L\subset L^1$ by
  Lemma~\ref{lem:daniell-ext-contains-source}, it suffices to show
  $L^1\subset L$.

  Let $f\in L^1$ and choose a series representation
  $f\overset{\varphi_n}{\simeq}\sum f_n$.
  Put
  \[
    Z:=\Bigl\{x\in X\;\Bigm|\;\sum_n\varphi_n(x)=+\infty\Bigr\} .
  \]
  From $|f_n|\le\varphi_n$ we get $\varphi_n\ge0$, and
  $\sum_nI(\varphi_n)<+\infty$, so $\{\varphi_n\}$ witnesses that $Z$ is an
  $I$-null set.

  Put $p_n:=f_n\lor0$ and $m_n:=(-f_n)\lor0$, so that $p_n,m_n\in L$ and
  \[
    f_n=p_n-m_n,
    \qquad
    0\le p_n,m_n\le|f_n|\le\varphi_n .
  \]
  Put
  \[
    P_N:=\sum_{n\le N}p_n\in L,
    \qquad
    M_N:=\sum_{n\le N}m_n\in L .
  \]
  Both are nondecreasing in $N$, and
  Lemma~\ref{lem:integration-preserve-order} gives
  \[
    I(P_N)\le\sum_{n\le N}I(\varphi_n)\le\sum_nI(\varphi_n)<+\infty ,
  \]
  and likewise for $M_N$.

  If $x\notin Z$ then $\sum_n\varphi_n(x)<+\infty$, so $\sum_np_n(x)$ and
  $\sum_nm_n(x)$ both converge. We therefore define
  \[
    P'(x):=
    \begin{cases}
      \sum_np_n(x) & (x\notin Z)\\
      0 & (x\in Z)
    \end{cases}
  \]
  and define $M'$ similarly. Both are real-valued functions on $X$.
  Let $P_N'$ be the function agreeing with $P_N$ on $X\setminus Z$ and vanishing
  on $Z$, and define $M_N'$ similarly.
  By hypothesis $(X,L,I)$ is closed under modification on null sets, so
  $P_N',M_N'\in L$.

  Since $P_N-P_N'$ vanishes on $X\setminus Z$,
  Lemma~\ref{lem:vanish-off-null-set} gives $I(P_N-P_N')=0$, that is,
  \[
    I(P_N')=I(P_N)\le\sum_nI(\varphi_n)<+\infty .
  \]
  Moreover $P_N'(x)\uparrow P'(x)$ for every $x\in X$ (for $x\notin Z$ by
  definition, and for $x\in Z$ both sides are $0$).
  By hypothesis $(X,L,I)$ satisfies the monotone convergence theorem, so
  $P'\in L$. Similarly $M'\in L$.

  Hence $P'-M'\in L$.
  If $x\notin Z$ then $\sum_nf_n(x)$ converges absolutely by
  Lemma~\ref{lem:series-exp-is-abs-conv}, so
  \[
    (P'-M')(x)=\sum_np_n(x)-\sum_nm_n(x)=\sum_nf_n(x)=f(x) .
  \]
  That is, $f$ agrees with the element $P'-M'$ of $L$ on $X\setminus Z$.
  Using closure under modification on null sets again, $f\in L$.
\end{proof}

\medskip
\noindent
\textbf{On the construction in (ii)$\Rightarrow$(i).}
The functions $P'$ and $M'$ built in the proof above are both monotonically
nondecreasing limits of elements of $L$ with bounded integrals, that is,
\textbf{upper functions} in the classical sense.
What the proof exhibits is therefore the representation of an $f\in L^1$ as
agreeing, off a null set, with the difference $P'-M'$ of two upper functions.
Constructing the Lebesgue integral as differences of upper functions is a
classical method, and a representation of this shape is well known in that
context.
We do not use this correspondence below, and so cite no particular source for
it.

\begin{corollary}\label{cor:mintbar-intrinsic}
  $\mathbf{DInt}$ is the full subcategory of $\mathbf{Int}$ consisting of the
  \textbf{integration systems that satisfy the monotone convergence theorem and
  are closed under modification on null sets}.
\end{corollary}
\begin{proof}
  By Theorem~\ref{thm:complete-intrinsic}.
\end{proof}

\begin{remark}\label{rem:not-about-mentioning-L1}
  The description above does not pass through the construction of the Daniell
  extension, but \textbf{that in itself is not where the point lies}.
  The condition $L^1=L$, when unfolded, also reads
  \begin{quote}
    every function admitting a series representation by elements of $L$ belongs
    to $L$,
  \end{quote}
  which uses nothing but $L$, $I$ and pointwise conditions.
  ``Whether $L^1$ is mentioned'' does not distinguish the two.

  What matters is that one unfamiliar closure condition splits into two
  conditions that occur independently in measure theory.
  We return to this point in
  Remark~\ref{rem:main-theorem-not-tautology}.
\end{remark}

\begin{example}
  Let us check this in the example of
  Remark~\ref{rem:D-I-not-left-adjoint} ($X=\{0,1,2\}$, $L=\{(a,b,b)\}$,
  $I(a,b,b)=a$).
  The function $g_n:=(0,1,1)\in L$ satisfies $g_n\ge0$ and $I(g_n)=0$, and
  $\sum_ng_n(x)=+\infty$ at $x=1,2$, so $Z=\{1,2\}$ is an $I$-null set.
  But $h=(0,1,0)$, obtained from $f=(0,0,0)\in L$ by modifying it on $Z$, does
  not belong to $L$.
  Hence this integration system is not closed under modification on null sets
  and, by Theorem~\ref{thm:complete-intrinsic}, is not Daniell closed.
  This agrees with the direct computation $L^1=\mathbb{R}^X\ne L$ in that
  example.

  On the other hand, in Example~\ref{ex:stone-fails} ($X=\{1,2\}$,
  $L=\{(a,2a)\}$, $I(a,2a)=a$), if $g_n=(c_n,2c_n)\ge0$ satisfies
  $\sum_nI(g_n)=\sum_nc_n<+\infty$ then $\sum_ng_n(x)<+\infty$ holds at both
  $x=1$ and $x=2$.
  Hence the only $I$-null set is the empty set, and the condition on
  modification on null sets holds vacuously.
  The monotone convergence theorem holds as well, so this integration system is
  Daniell closed by Theorem~\ref{thm:complete-intrinsic}, in agreement with
  $L^1=L$ as stated in that example.
\end{example}

\begin{remark}[neither condition can be dropped]\label{rem:two-conditions-independent}
  The two conditions in (ii) of Theorem~\ref{thm:complete-intrinsic} are
  independent.

  \textbf{That the monotone convergence theorem cannot be dropped} is shown by
  Example~\ref{ex:counting} ($X=\mathbb{N}$, the functions vanishing outside
  finitely many $n$, and summation).
  In this integration system, $g\in L$ with $g\ge0$ satisfies $g(x)\le I(g)$ for
  every $x$, so for a nonnegative sequence $\{g_n\}\subset L$ with
  $\sum_nI(g_n)<+\infty$ we have $\sum_ng_n(x)<+\infty$ at every $x$.
  Hence the only $I$-null set is the empty set, and the condition on
  modification on null sets holds vacuously.
  On the other hand
  \[
    f_N:=\sum_{k\le N}2^{-k}1_{\{k\}}\in L
  \]
  satisfies $I(f_N)\le1$ while increasing monotonically pointwise to
  $f(k)=2^{-k}$, and $f\notin L$ because $f$ does not have finite support.
  That is, the monotone convergence theorem fails.
  Indeed this integration system is not Daniell closed, and $L^1=\ell^1$.

  \textbf{That modification on null sets cannot be dropped} is shown by the
  example of Remark~\ref{rem:D-I-not-left-adjoint}.
  The integration system of that example satisfies the monotone convergence
  theorem but, as seen in the example above, is not closed under modification on
  null sets, hence not Daniell closed.
\end{remark}

\begin{remark}
  The two conditions in (ii) of Theorem~\ref{thm:complete-intrinsic} correspond
  to two properties that $L^1(\mu)$ satisfies in measure theory.
  The monotone convergence theorem corresponds directly, and closure under
  modification on null sets corresponds to the fact that $L^1(\mu)$ contains the
  functions measurable with respect to the completion of $\mu$.
  In Proposition~\ref{prop:l1-decomposition} of
  Subsection~\ref{sec:why-complete} we will see that the gap between the two
  reflections amounts exactly to the functions supported on null sets; the
  present theorem explains the reason for that from the side of Daniell
  closedness.
  That is, \textbf{to be an object of $\mathbf{DInt}$ one must be closed under
  modification on null sets}.
\end{remark}

\begin{remark}[the decomposition collapses in the extended-real formalism]\label{rem:extended-real-collapse}
  That the two conditions of Theorem~\ref{thm:complete-intrinsic} are
  independent \textbf{depends on this paper's restriction of functions to real
  values}.
  If the range is widened to
  $\overline{\mathbb{R}}=[-\infty,+\infty]$, then condition (b) follows from
  condition (a).

  Indeed \cite{CWS} adopts that formalism.
  It defines the counterpart of a Riesz space inside
  $\overline{\mathbb{R}}^X$ under the name Riesz lattice, and calls the triple
  corresponding to our integration system a Daniell space.
  It then gives the two conditions separate names.
  \begin{itemize}
    \item Its Definition 1.4.1 calls a triple \textbf{closed} if, for every
      monotone sequence whose integrals are bounded, the pointwise limit again
      belongs to the space and the integral commutes with the limit.
      This is our Definition~\ref{def:satisfies-mct} (satisfying the monotone
      convergence theorem).
    \item Its Definition 1.4.7 calls a triple \textbf{complete} if
      $f\le h\le g$ together with $\ell(f)=\ell(g)$ implies that $h$ belongs to
      the space as well.
      This sandwich condition is in substance our
      Definition~\ref{def:mod-closed}.
  \end{itemize}
  Its Corollary 1.4.8, however, states that
  \begin{quote}
    every closed Daniell space is complete.
  \end{quote}
  That is, \textbf{in the extended-real formalism (a)$\Rightarrow$(b), and the
  two conditions are not independent.}

  The reason is the following. Let $Z$ be an $I$-null set and take the sequence
  $\{g_n\}$ of Definition~\ref{def:null-set}.
  The partial sums $\sum_{n\le N}g_n$ are monotonically nondecreasing with
  bounded integrals, so under (a) their pointwise limit belongs to the space.
  If extended-real values are allowed, that limit is an element of the space as
  a \emph{function} taking the value $+\infty$ on $Z$.
  Null sets can then be named from inside the space, as the sets where an
  element diverges to $\infty$, and (b) follows from the axioms of a Riesz
  lattice, which originate in the fact that $\infty-\infty$ is undefined.

  \textbf{This argument does not go through here.}
  Since $L\subset\mathbb{R}^X$, the limit above cannot be an element of $L$.
  A null set cannot be captured from inside $L$ as divergence to $\infty$, and
  it is precisely for this reason that
  Definition~\ref{def:null-set} takes the form of a series and that condition
  (b) remains as an independent condition.
  The counterexample of
  Remark~\ref{rem:two-conditions-independent} ($X=\{0,1,2\}$,
  $L=\{(a,b,b)\}$, $I(a,b,b)=a$) was an integration system satisfying the
  monotone convergence theorem but not closed under modification on null sets.
  That is, it is also a \textbf{counterexample to the real-valued form of
  Corollary 1.4.8}.

  This is not a question of which formalism is better.
  The correspondence is that \textbf{widening the range makes null sets nameable
  inside the space, at the cost of collapsing the two closure conditions into
  one}.
  Theorem~\ref{thm:complete-intrinsic} is a decomposition visible only in a
  formalism restricted to real values.
\end{remark}

\subsection{The $L^1$ seminorm and approximation}

The content of this subsection is also known, and corresponds to \S4 and \S6 of
\cite{BM}.
We give proofs because it is used essentially in
Section~\ref{chap:operator}.

\begin{definition}\label{def:l1-seminorm}
  Let $(X,L,I)$ be an integration system on a set $X$.
  For each $f\in L$ we put
  \[
    \|f\|_I:=I(|f|) .
  \]
  Since $L$ is a Riesz space we have $|f|\in L$, so this value is defined.
\end{definition}

\begin{lemma}\label{lem:l1-seminorm}
  $\|\cdot\|_I$ is a seminorm on $L$. That is, for all $f,g\in L$ and
  $a\in\mathbb{R}$,
  \[
    \|af\|_I=|a|\,\|f\|_I,
    \qquad
    \|f+g\|_I\le\|f\|_I+\|g\|_I.
  \]
\end{lemma}
\begin{proof}
  The first identity follows from $|af|=|a||f|$ and linearity of $I$.
  Moreover $|f+g|\le|f|+|g|$ and
  Lemma~\ref{lem:integration-preserve-order} give
  $I(|f+g|)\le I(|f|)+I(|g|)$.
\end{proof}

\begin{remark}
  In general $\|f\|_I=0$ does not imply $f=0$.
  The operation that collapses this discrepancy is treated in
  Subsection~\ref{sec:ae}.
\end{remark}

\begin{lemma}[approximation of integrable functions by $L$]\label{lem:l1-approx-by-source}
  Let $(X,L,I)$ be an integration system on a set $X$ and let
  $(X,L^1,I^1)$ be its Daniell extension.
  Then for every $f\in L^1$ there is a sequence of functions
  $\{s_N\}\subset L$ satisfying
  \[
    \|f-s_N\|_{I^1}=I^1(|f-s_N|)\longrightarrow0
    \qquad(N\to\infty) .
  \]
  That is, $L$ is dense in $L^1$ with respect to $\|\cdot\|_{I^1}$.
\end{lemma}
\begin{proof}
  Choose a series representation of $f$, say
  \[
    f\overset{\varphi_n}{\simeq}\sum_{n=1}^{\infty}f_n ,
  \]
  and put
  \[
    s_N:=\sum_{n=1}^Nf_n\in L,
    \qquad
    r_N:=f-s_N .
  \]
  Since finitely many values $\varphi_n(x)$ are always finite,
  \[
    r_N\overset{\varphi_{N+m}}{\simeq}
    \sum_{m=1}^{\infty}f_{N+m}
  \]
  is a series representation of $r_N$.

  For each $m\ge1$ put
  \[
    g_{N,m}:=
    \left|\sum_{k=1}^{m}f_{N+k}\right|
    -\left|\sum_{k=1}^{m-1}f_{N+k}\right| ,
  \]
  where the empty sum is $0$. Since $L$ is a Riesz space we have
  $g_{N,m}\in L$, and the triangle inequality for absolute values gives
  \[
    |g_{N,m}|\le |f_{N+m}|\le\varphi_{N+m} .
  \]
  Moreover, if $\sum_m\varphi_{N+m}(x)<+\infty$ then, by telescoping the sum and
  continuity of the absolute value,
  \[
    \sum_{m=1}^{\infty}g_{N,m}(x)
    =\left|\sum_{m=1}^{\infty}f_{N+m}(x)\right|
    =|r_N(x)| .
  \]
  Hence
  \[
    |r_N|\overset{\varphi_{N+m}}{\simeq}
    \sum_{m=1}^{\infty}g_{N,m}
  \]
  is a series representation of $|r_N|$. By positivity of the integral,
  \[
    \begin{aligned}
      0\le I^1(|f-s_N|)
      &=\sum_{m=1}^{\infty}I(g_{N,m})\\
      &\le\sum_{m=1}^{\infty}I(|g_{N,m}|)\\
      &\le\sum_{m=1}^{\infty}I(\varphi_{N+m})
      \longrightarrow0 ,
    \end{aligned}
  \]
  which is the assertion.
\end{proof}

\begin{corollary}\label{cor:tail-estimate}
  Let $(X,L,I)$ be an integration system and suppose $\{h_j\}\subset L$
  satisfies $\sum_j I(|h_j|)<+\infty$.
  Then the function
  \[
    h(x):=\begin{cases}
      \sum_j h_j(x) & \text{if $\sum_j|h_j(x)|<+\infty$}\\
      0 & \text{otherwise}
    \end{cases}
  \]
  belongs to $L^1$, and for each $J\ge1$
  \[
    \left\|h-\sum_{j=1}^{J}h_j\right\|_{I^1}\le\sum_{j>J}I(|h_j|)
  \]
  holds.
\end{corollary}
\begin{proof}
  Putting $\varphi_j:=|h_j|\in L$, the pair $\{h_j\},\{\varphi_j\}$ is a series
  representation of $h$, so $h\in L^1$.
  Moreover $h-\sum_{j\le J}h_j$ has $\{h_{J+m}\}_{m\ge1}$ and
  $\{\varphi_{J+m}\}_{m\ge1}$ as a series representation, so the same estimate
  as at the end of the proof of Lemma~\ref{lem:l1-approx-by-source} gives the
  stated inequality.
\end{proof}

\begin{theorem}[$\|\cdot\|$ completeness]\label{thm:l1-norm-complete}
  Let $(X,L,I)$ be a Daniell closed integration system.
  Then $(L,\|\cdot\|_I)$ is complete.
  That is, for every Cauchy sequence $\{g_k\}\subset L$ with respect to
  $\|\cdot\|_I$ there is a $g\in L$ with $\|g-g_k\|_I\to0$.
\end{theorem}
\begin{proof}
  Let $\{g_k\}$ be a Cauchy sequence. Choose a subsequence $\{g_{k_j}\}$ so that
  \[
    \|g_{k_{j+1}}-g_{k_j}\|_I<2^{-j} ,
  \]
  and put
  \[
    h_1:=g_{k_1},\qquad h_{j+1}:=g_{k_{j+1}}-g_{k_j} .
  \]
  Then $\sum_j I(|h_j|)<+\infty$.
  By Corollary~\ref{cor:tail-estimate} there is a corresponding $h\in L^1$ with
  \[
    \left\|h-g_{k_J}\right\|_{I^1}
    =\left\|h-\sum_{j=1}^{J}h_j\right\|_{I^1}
    \le\sum_{j>J}I(|h_j|)
    \longrightarrow0 .
  \]
  Since $(X,L,I)$ is Daniell closed we have $L^1=L$, and
  $I^1|_L=I$ by Lemma~\ref{lem:daniell-ext-contains-source}.
  Hence $h\in L$ and $\|h-g_{k_J}\|_I\to0$.
  A Cauchy sequence with a convergent subsequence converges to the same limit as
  a whole, so $\|h-g_k\|_I\to0$.
\end{proof}

\subsection{Null functions and equality almost everywhere}\label{sec:ae}

The seminorm $\|\cdot\|_I$ is in general not a norm.
We set up here the operation that collapses its kernel.
This quotient is used essentially in Section~\ref{chap:operator}.

\subsubsection{Null functions and the quotient}

\begin{definition}\label{def:null-function}
  Let $(X,L,I)$ be an integration system on a set $X$. We put
  \[
    \mathcal{N}_I:=\{f\in L\mid I(|f|)=0\}
  \]
  and call the functions belonging to $\mathcal{N}_I$ the \textbf{null functions}
  of $(X,L,I)$.
\end{definition}

\begin{lemma}\label{lem:null-is-ideal}
  $\mathcal{N}_I$ is an ideal of $L$.
  That is, $\mathcal{N}_I$ is a linear subspace, and if $f\in\mathcal{N}_I$ and
  $f'\in L$ satisfies $|f'|\le|f|$ then $f'\in\mathcal{N}_I$.
\end{lemma}
\begin{proof}
  $\mathcal{N}_I$ is the kernel of the seminorm $\|\cdot\|_I$, hence a linear
  subspace (Lemma~\ref{lem:l1-seminorm}).
  Moreover, if $|f'|\le|f|$ then
  Lemma~\ref{lem:integration-preserve-order} gives $I(|f'|)\le I(|f|)=0$, and
  positivity of $I$ gives $I(|f'|)=0$.
\end{proof}

\begin{definition}\label{def:null-quotient}
  We write the quotient space as
  \[
    \widetilde{L}:=L/\mathcal{N}_I
  \]
  and denote the quotient map by $\pi_I:L\to\widetilde{L}$.
\end{definition}

\begin{lemma}\label{lem:quotient-riesz}
  $\widetilde{L}$ carries the structure of a Riesz space for the order defined by
  declaring $\pi_I(f)\ge0$ to mean ``there is an $f'\in\mathcal{N}_I$ with
  $f+f'\ge0$'', and $\pi_I$ is a Riesz homomorphism.
  That is,
  \[
    \pi_I(f\vee g)=\pi_I(f)\vee\pi_I(g),
    \qquad
    \pi_I(|f|)=|\pi_I(f)| .
  \]
  Furthermore
  \[
    \widetilde{I}:\widetilde{L}\to\mathbb{R};\quad \pi_I(f)\mapsto I(f)
  \]
  is a well-defined positive linear functional, and
  $\|\pi_I(f)\|:=\widetilde{I}(|\pi_I(f)|)=\|f\|_I$ is a norm.
\end{lemma}
\begin{proof}
  That the quotient of a Riesz space by an ideal is a Riesz space and that the
  quotient map is a Riesz homomorphism is general theory (see for instance
  \cite{LuxemburgZaanen}).
  By Lemma~\ref{lem:null-is-ideal} this applies here.

  $\widetilde{I}$ is well defined because $f\in\mathcal{N}_I$ implies
  $|I(f)|\le I(|f|)=0$
  (Lemma~\ref{lem:triangle-ineq-for-integration}).
  Positivity and linearity follow from those of $I$.
  Finally
  $\widetilde{I}(|\pi_I(f)|)=\widetilde{I}(\pi_I(|f|))=I(|f|)=\|f\|_I$, and this
  vanishes only when $f\in\mathcal{N}_I$, that is, when $\pi_I(f)=0$.
\end{proof}

\begin{lemma}\label{lem:quotient-is-banach}
  If $(X,L,I)$ is a Daniell closed integration system then
  $(\widetilde{L},\|\cdot\|)$ is a Banach space.
\end{lemma}
\begin{proof}
  If $\{\pi_I(f_k)\}$ is Cauchy then $\{f_k\}$ is a Cauchy sequence with respect
  to $\|\cdot\|_I$.
  By Theorem~\ref{thm:l1-norm-complete} its limit $f\in L$ exists, and
  $\|\pi_I(f)-\pi_I(f_k)\|=\|f-f_k\|_I\to0$.
\end{proof}

\begin{lemma}[compatibility of the Daniell extension with the quotient]\label{lem:quotient-embedding}
  Let $(X,L,I)$ be an integration system on a set $X$ and let
  $(X,L^1,I^1)$ be its Daniell extension.
  Then
  \[
    \mathcal{N}_{I^1}\cap L=\mathcal{N}_I .
  \]
  Hence the inclusion $L\subset L^1$ induces an injective Riesz homomorphism
  \[
    j_I:\widetilde{L}\longrightarrow\widetilde{L^1}
  \]
  satisfying
  \[
    \pi_{I^1}|_L=j_I\circ\pi_I .
  \]
  Furthermore $j_I$ preserves the norm.
\end{lemma}
\begin{proof}
  By Lemma~\ref{lem:daniell-ext-contains-source} we have $|f|\in L$ and
  $I^1(|f|)=I(|f|)$ for every $f\in L$.
  Hence $I^1(|f|)=0$ and $I(|f|)=0$ are equivalent, and we obtain
  $\mathcal{N}_{I^1}\cap L=\mathcal{N}_I$.

  It follows that the kernel of $\pi_{I^1}|_L:L\to\widetilde{L^1}$ is
  $\mathcal{N}_I$, so an injective linear map
  $j_I:\widetilde{L}\to\widetilde{L^1}$ is induced, satisfying
  $\pi_{I^1}|_L=j_I\circ\pi_I$.
  Both $\pi_I$ and $\pi_{I^1}$ are Riesz homomorphisms
  (Lemma~\ref{lem:quotient-riesz}) and $\pi_I$ is surjective, so $j_I$ is a
  Riesz homomorphism as well.
  Finally, for $f\in L$,
  \[
    \|j_I(\pi_I(f))\|=\|\pi_{I^1}(f)\|=I^1(|f|)=I(|f|)=\|\pi_I(f)\| .
  \]
\end{proof}

\subsubsection{ae-equivalence of morphisms}

\begin{definition}\label{def:ae-equal}
  Let $(X,L,I)$ and $(Y,M,J)$ be integration systems on sets $X$ and
  $Y$ respectively.
  We say that two morphisms of integration systems
  \[
    F,G:(X,L,I)\to(Y,M,J)
  \]
  are \textbf{equal almost everywhere} if
  \[
    J\bigl(|F^*f-G^*f|\bigr)=0
  \]
  holds for every $f\in L$, that is, if $F^*f-G^*f\in\mathcal{N}_J$ holds for
  every $f\in L$.
  In this case we write
  \[
    F=_{\mathrm{ae}}G .
  \]
\end{definition}

\begin{lemma}\label{lem:ae-is-congruence}
  For the morphisms of integration systems between each pair of objects, the
  relation $=_{\mathrm{ae}}$ is an equivalence relation and is compatible with
  composition of morphisms.
  That is, if morphisms of integration systems
  \[
    F,F':(X,L,I)\to(Y,M,J),
    \qquad
    G,G':(Y,M,J)\to(Z,N,K)
  \]
  satisfy
  \[
    F=_{\mathrm{ae}}F',
    \qquad
    G=_{\mathrm{ae}}G' ,
  \]
  then
  \[
    G\circ F=_{\mathrm{ae}}G'\circ F' .
  \]
\end{lemma}
\begin{proof}
  Reflexivity and symmetry are immediate from the definition. As for
  transitivity, suppose morphisms of integration systems
  $F,F',F'':(X,L,I)\to(Y,M,J)$ satisfy $F=_{\mathrm{ae}}F'$ and
  $F'=_{\mathrm{ae}}F''$. For every $f\in L$,
  \[
    \begin{aligned}
      0
      &\le J\bigl(|F^*f-(F'')^*f|\bigr)\\
      &\le J\bigl(|F^*f-(F')^*f|\bigr)
        +J\bigl(|(F')^*f-(F'')^*f|\bigr)
      =0 ,
    \end{aligned}
  \]
  so $F=_{\mathrm{ae}}F''$.

  Next we show compatibility with composition. For every $f\in L$, since
  morphisms of Riesz systems preserve absolute values and morphisms of
  integration systems preserve the integral,
  \[
    \begin{aligned}
      &K\bigl(|G^*F^*f-(G')^*(F')^*f|\bigr)\\
      &\quad\le
      K\bigl(|G^*F^*f-G^*(F')^*f|\bigr)
      +K\bigl(|G^*(F')^*f-(G')^*(F')^*f|\bigr)\\
      &\quad=
      J\bigl(|F^*f-(F')^*f|\bigr)
      +K\bigl(|G^*(F')^*f-(G')^*(F')^*f|\bigr)
      =0.
    \end{aligned}
  \]
  Therefore $G\circ F=_{\mathrm{ae}}G'\circ F'$.
\end{proof}

\begin{definition}\label{def:int-ae}
  The category having the same objects as the category $\mathbf{Int}$ of
  integration systems and whose morphisms are the equivalence classes under
  $=_{\mathrm{ae}}$ is written
  \[
    \mathbf{Int}_{\mathrm{ae}}
  \]
  and called the \textbf{ae-category of integration systems}.
  By Lemma~\ref{lem:ae-is-congruence}, composition of morphisms is defined
  independently of the choice of representatives.

  Among the objects of $\mathbf{Int}_{\mathrm{ae}}$, we write the full
  subcategory consisting of the integration systems satisfying the monotone
  convergence theorem as
  \[
    \mathbf{MInt}_{\mathrm{ae}},
  \]
  and the full subcategory consisting of the Daniell closed integration systems
  as
  \[
    \mathbf{DInt}_{\mathrm{ae}} .
  \]
\end{definition}

\subsection{Relation to measures}\label{sec:measure}

Up to here this paper has treated integration without using measures.
In this subsection we state the correspondence between integration systems and
measure theory.
The main results of this paper do not use this correspondence, but we record that
the notion of an integration system does recover measure theory.
Everything below is known, and we refer to the literature for proofs.

\subsubsection{From measures to integration systems}

This direction is the easy one.

\begin{proposition}\label{prop:measure-gives-integral-system}
  Let $(X,\Sigma,\mu)$ be a measure space, let $L$ be the set of all
  $\Sigma$-measurable simple functions that are integrable with respect to
  $\mu$, and let $I(f):=\int_Xf\,d\mu$.
  Then $(X,L,I)$ is an integration system.
\end{proposition}

\noindent
Positivity is immediate, and verifying the continuity axiom uses
$\sigma$-additivity of $\mu$.
For a proof see \cite{BM}, \S7.
The set of step functions in Example~\ref{ex:step-functions} is a subsystem of
the set of all simple functions with respect to Lebesgue measure.
Finite unions of half-open intervals do not form a $\sigma$-algebra, so it is not
itself a special case of this proposition.

\subsubsection{The Stone condition}

The converse direction, that of extracting a measure from an integration system,
needs one condition.

\begin{definition}\label{def:stone-condition}
  We say that an integration system $(X,L,I)$ on a set $X$ satisfies
  the \textbf{Stone condition} if
  \[
    f\wedge1\in L
  \]
  holds for every $f\in L$. Here $1$ is the constant function $1$ on $X$.
\end{definition}

\noindent
Note that $L$ itself need not contain $1$.
What is required is that the function obtained by truncating $f$ at $1$ belong to
$L$ again.
This condition is the standard one, and a Riesz space satisfying it is called a
\textbf{Stone vector lattice}.

This condition is needed because we want to define the measure of a set $A$ with
$1_A\in L^1$ by
\[
  \mu(A):=I^1(1_A) .
\]
If $L^1$ contains no indicator function at all, this definition cannot even be
started.

\begin{remark}\label{rem:mu-only-on-finite}
  The formula above is meaningful only for sets with $1_A\in L^1$.
  It cannot be used for general measurable sets, and in particular not for sets
  of infinite measure.

  Looking at Example~\ref{ex:counting} ($X=\mathbb{N}$, the functions of finite
  support, and summation), $L^1$ is the set $\ell^1$ of all absolutely summable
  sequences and the measure one should obtain is the counting measure.
  But $1_{\mathbb{N}}\notin\ell^1$, so $I^1(1_{\mathbb{N}})$ is undefined, while
  $\mu(\mathbb{N})=\infty$.

  To construct a general measure one must therefore start from the sets with
  $1_A\in L^1$ and pass through an upper integral or an outer measure.
  Theorem~\ref{thm:daniell-stone} below states only the result of that
  construction.
\end{remark}

\begin{example}\label{ex:stone-holds}
  Both Example~\ref{ex:counting} and Example~\ref{ex:step-functions} satisfy the
  Stone condition.
  Indeed, truncating a function of finite support at $1$ again has finite
  support, and truncating a step function at $1$ is again a step function.
\end{example}

\begin{example}[an integration system not satisfying the Stone condition]\label{ex:stone-fails}
  Let $X=\{1,2\}$ and put
  \[
    L:=\{(a,2a)\mid a\in\mathbb{R}\}\subset\mathbb{R}^X,
    \qquad
    I(a,2a):=a .
  \]
  Then $L$ is the one-dimensional subspace spanned by $(1,2)$, and
  $|(a,2a)|=|a|\cdot(1,2)\in L$, so $L$ is a Riesz space on $X$.
  Moreover $I$ is clearly positive linear, and continuity holds because $X$ is
  finite.
  Hence $(X,L,I)$ is an integration system.

  But $f=(1,2)$ gives $f\wedge1=(1,1)\notin L$, so the Stone condition fails.

  In this integration system $L^1=L$.
  Indeed, the dominating functions occurring in a series representation have the
  form $\varphi_n=(c_n,2c_n)$ with $c_n\ge0$, and if
  $\sum_nI(\varphi_n)=\sum_nc_n<+\infty$ then $\sum_n\varphi_n(x)<+\infty$ holds
  at both $x=1$ and $x=2$.
  Hence condition (3) of a series representation is imposed on all of $X$, and
  $f$ is confined to being the sum of an absolutely convergent series of
  elements of $L$, that is, an element of $L$.

  Consequently $1_{\{1\}}=(1,0)\notin L^1$, and $\mu(\{1\}):=I^1(1_{\{1\}})$
  cannot be defined.
\end{example}

\subsubsection{The Daniell--Stone representation theorem}

\begin{theorem}[Daniell--Stone]\label{thm:daniell-stone}
  Let $(X,L,I)$ be an integration system satisfying the Stone condition.
  Then there exist a $\sigma$-algebra $\Sigma$ on $X$ and a measure $\mu$ such
  that every $f\in L$ is $\mu$-integrable and
  \[
    I(f)=\int_Xf\,d\mu .
  \]
\end{theorem}

\noindent
For a proof see \cite{Stone}, \cite{Royden} Chapter 16, or \cite{CFW}.
Chapter 16 of \cite{Royden} is entitled ``The Daniell Integral''; its \S2 is the
extension theorem, \S3 uniqueness, and \S4 measurability and measure.
Here $\Sigma$ is chosen so that the functions belonging to $L^1$ are measurable,
and $\mu$ is constructed through an outer measure, as in
Remark~\ref{rem:mu-only-on-finite}, so as to satisfy $\mu(A)=I^1(1_A)$ for the
sets $A$ with $1_A\in L^1$.

\begin{example}
  Applying this theorem to Example~\ref{ex:step-functions}, the measure obtained
  is Lebesgue measure on $[0,1)$ and the Daniell extension $L^1$ coincides with
  the set of all Lebesgue integrable functions.
  This is how the fact stated in Example~\ref{ex:lebesgue} looks from the side
  of measure theory.
\end{example}

\begin{remark}
  This paper does not impose the Stone condition.
  The results of Section~\ref{chap:functoriality} and
  Section~\ref{chap:operator} hold without it.
  Indeed the integration system of Example~\ref{ex:stone-fails} is an object of
  $\mathbf{Int}$ as well, and is moreover Daniell closed.

  One cannot, however, conclude from this that ``the framework of this paper is
  wider than measure theory''.
  The integral of Example~\ref{ex:stone-fails} can be written as
  $I(a,2a)=\int_Xf\,d\mu$ by taking $\mu(\{1\})=1$ and $\mu(\{2\})=0$
  (and such a $\mu$ is moreover not unique: any $\mu$ with
  $\mu(\{1\})+2\mu(\{2\})=1$ will do).
  What that example shows is that without the Stone condition one
  \textbf{cannot canonically recover a measure from inside the Daniell extension
  using indicator functions}, and that \textbf{the representing measure is not
  uniquely determined}.

  Put precisely, this paper does not assume the Stone condition and does not take
  representation by a measure as a premise of the theory.
\end{remark}

\begin{problem}\label{prob:stone-measure-spaces}
  Is every Daniell closed integration system satisfying the Stone condition
  realized as the $L^1$ of a measure space, and does this correspondence lift to
  an equivalence of categories?
  That is, is the full subcategory of $\mathbf{DInt}$ consisting of the objects
  that satisfy the Stone condition equivalent to a suitable category of measure
  spaces and measure-preserving maps?
\end{problem}

\noindent
This question was raised by T.~Fritz in personal communication (2026).
It does not hold, however, in the form of an equivalence with the category of
measurable spaces.
A morphism of integration systems preserves the integral, $J(F^*f)=I(f)$, so the
integral itself is data that distinguishes objects up to isomorphism.
Indeed, let $X=\{*\}$, $L=\mathbb{R}$ and $I_c(a):=ca$ for $c>0$.
Since $L^1\subset\mathbb{R}^X=\mathbb{R}=L$, each of these is a Daniell closed
integration system satisfying the Stone condition.
But $(X,\mathbb{R},I_c)\not\cong(X,\mathbb{R},I_{c'})$ for $c\ne c'$, while the
underlying measurable space is a one-point space in every case.
What should be asked about is therefore measure spaces, not measurable spaces.
Moreover, by Theorem~\ref{thm:complete-intrinsic} a Daniell closed integration
system is closed under modification on null sets, so the measures to be matched
are expected to be the complete ones.

